\documentclass[11pt,a4paper]{article}

% ---------------------------------------------------------------
% Tipografia: Palatino, versalete, uma cor de destaque (bord\^o)
% Mesmo pre\^ambulo da lista de termologia.
% ---------------------------------------------------------------
\usepackage[utf8]{inputenc}
\usepackage[T1]{fontenc}
\usepackage[brazil]{babel}
\usepackage[sc,osf]{mathpazo}
\usepackage{microtype}
\usepackage[top=2.4cm, bottom=2.6cm, left=2.6cm, right=2.6cm]{geometry}
\usepackage{amsmath}
\usepackage{amssymb}
\usepackage{xcolor}
\usepackage{tcolorbox}
\tcbuselibrary{skins, breakable}
\usepackage{enumitem}
\usepackage{booktabs}
\usepackage{array}
\usepackage{needspace}
\usepackage{hyperref}
\usepackage{graphicx}
\graphicspath{{figuras/}}

\definecolor{vinho}{RGB}{122, 27, 42}
\definecolor{grafite}{RGB}{70, 70, 70}
\hypersetup{colorlinks=true, linkcolor=vinho, urlcolor=vinho}

\linespread{1.06}
\setlength{\parindent}{0pt}
\setlength{\parskip}{0.45em}

% ---------------------------------------------------------------
% Formata\c{c}\~ao das quest\~oes
% ---------------------------------------------------------------
\newcommand{\fonte}[1]{{\small\scshape\color{grafite}#1}}

\newenvironment{questao}[2]{%
  \par\addvspace{1.6em}
  \needspace{12\baselineskip}
  {\large\bfseries\color{vinho}Quest\~ao #1}\hfill #2\\[-0.55em]
  {\color{black!25}\rule{\linewidth}{0.4pt}}
  \par\nobreak\smallskip
  % a quest\~ao inteira numa caixa: nunca quebra entre p\'aginas
  \noindent\begin{minipage}{\linewidth}\setlength{\parskip}{0.45em}%
}{\end{minipage}\par\addvspace{0.4em}}

\newcommand{\parte}[1]{%
  \needspace{18\baselineskip}
  \vspace{1.8em}
  {\Large\bfseries #1}\\[-0.55em]
  {\color{black!30}\rule{\linewidth}{0.5pt}}
  \par\nobreak\smallskip
}

\newtcolorbox{painel}{blanker, breakable, left=12pt,
  borderline west={1.4pt}{0pt}{black!30},
  before skip=11pt, after skip=11pt}

% figura recortada do caderno oficial do INEP
\newcommand{\figuraoriginal}[2]{%
  \begin{center}
    \includegraphics[width=#1]{#2}
  \end{center}}

% refer\^encia bibliogr\'afica do enunciado, como na prova
\newcommand{\fontebib}[1]{{\raggedleft\scriptsize\color{grafite}#1\par}}

\setlist[enumerate,1]{label=\textbf{\Alph*)}, leftmargin=2em, nosep}

% ---------------------------------------------------------------
\begin{document}
% ---------------------------------------------------------------

\thispagestyle{empty}

{\scshape\color{vinho}cursinho solid\'ario \,\textperiodcentered\, f\'isica}\\[-0.5em]
{\color{vinho}\rule{\linewidth}{1.6pt}}\\[0.9em]
{\huge\bfseries Simulado ENEM}\\[0.35em]
{\LARGE Ondas, Calor, Energia e Newton}\\[0.7em]
\hfill {\color{grafite}2026}\\[-0.3em]
{\color{black!30}\rule{\linewidth}{0.5pt}}

\vspace{0.6em}


% ===================================================================
\parte{Ondulat\'oria}
% ===================================================================

\begin{questao}{1}{\fonte{ENEM 2025 \,\textperiodcentered\, caderno amarelo \,\textperiodcentered\, Q128}}
\textbf{Carros el\'etricos p\~oem fim ao longo reinado da r\'adio AM}

As montadoras de autom\'oveis est\~ao em uma corrida rumo a um futuro
eletrificado, e r\'adios AM est\~ao ficando pelo caminho. O problema, dizem
os especialistas, \'e que os motores dos ve\'iculos el\'etricos geram
frequ\^encias eletromagn\'eticas do mesmo comprimento de onda que os sinais
das r\'adios AM. ``Voc\^e tem dois sinais que colidem um com o outro e se
cancelam'', disse o diretor de inova\c{c}\~ao de motores e tecnologia de uma
fabricante de autope\c{c}as.

\fontebib{Dispon\'ivel em: www1.folha.uol.com.br. Acesso em: 20 nov. 2018 (adaptado).}

O problema mencionado refere-se a qual fen\^omeno ondulat\'orio?
\begin{enumerate}
  \item Difra\c{c}\~ao.
  \item Reflex\~ao.
  \item Refra\c{c}\~ao.
  \item Resson\^ancia.
  \item Interfer\^encia.
\end{enumerate}
\end{questao}

\begin{questao}{2}{\fonte{ENEM 2013 \,\textperiodcentered\, caderno azul \,\textperiodcentered\, Q65}}
Uma manifesta\c{c}\~ao comum das torcidas em est\'adios de futebol \'e a ola
mexicana. Os espectadores de uma linha, sem sair do lugar e sem se
deslocarem lateralmente, ficam de p\'e e se sentam, sincronizados com os da
linha adjacente. O efeito coletivo se propaga pelos espectadores do
est\'adio, formando uma onda progressiva.

Calcula-se que a velocidade de propaga\c{c}\~ao dessa ``onda humana'' \'e
45 km/h, e que cada per\'iodo de oscila\c{c}\~ao cont\'em 16 pessoas, que se
levantam e sentam organizadamente e distanciadas entre si por 80 cm.

Nessa ola mexicana, a frequ\^encia da onda, em hertz, \'e um valor mais
pr\'oximo de
\begin{enumerate}
  \item 0,3.
  \item 0,5.
  \item 1,0.
  \item 1,9.
  \item 3,7.
\end{enumerate}
\end{questao}

\begin{questao}{3}{\fonte{ENEM 2020 \,\textperiodcentered\, caderno amarelo \,\textperiodcentered\, Q109}}
Dois engenheiros est\~ao verificando se uma cavidade perfurada no solo est\'a
de acordo com o planejamento de uma obra, cuja profundidade requerida \'e de
30 m. O teste \'e feito por um dispositivo denominado oscilador de \'audio de
frequ\^encia vari\'avel, que permite relacionar a profundidade com os valores
da frequ\^encia de duas resson\^ancias consecutivas, assim como em um tubo
sonoro fechado. A menor frequ\^encia de resson\^ancia que o aparelho mediu
foi 135 Hz. Considere que a velocidade do som dentro da cavidade perfurada
\'e de 360 m\,s$^{-1}$.

Se a profundidade estiver de acordo com o projeto, qual ser\'a o valor da
pr\'oxima frequ\^encia de resson\^ancia que ser\'a medida?
\begin{enumerate}
  \item 137 Hz.
  \item 138 Hz.
  \item 141 Hz.
  \item 144 Hz.
  \item 159 Hz.
\end{enumerate}
\end{questao}


% ===================================================================
\parte{Calorimetria}
% ===================================================================

\begin{questao}{4}{\fonte{ENEM 2023 \,\textperiodcentered\, caderno amarelo \,\textperiodcentered\, Q95}}
Em uma ind\'ustria aliment\'icia, para produ\c{c}\~ao de doce de leite,
utiliza-se um tacho de parede oca com uma entrada para vapor de \'agua a
120 $^\circ$C e uma sa\'ida para \'agua l\'iquida em equil\'ibrio com o vapor
a 100 $^\circ$C. Ao passar pela parte oca do tacho, o vapor de \'agua
transforma-se em l\'iquido, liberando energia. A parede transfere essa
energia para o interior do tacho, resultando na evapora\c{c}\~ao de \'agua e
consequente concentra\c{c}\~ao do produto.

No processo de concentra\c{c}\~ao do produto, \'e utilizada energia proveniente
\begin{enumerate}
  \item somente do calor latente de vaporiza\c{c}\~ao.
  \item somente do calor latente de condensa\c{c}\~ao.
  \item do calor sens\'ivel e do calor latente de vaporiza\c{c}\~ao.
  \item do calor sens\'ivel e do calor latente de condensa\c{c}\~ao.
  \item do calor latente de condensa\c{c}\~ao e do calor latente de
        vaporiza\c{c}\~ao.
\end{enumerate}
\end{questao}

\begin{questao}{5}{\fonte{ENEM 2020 \,\textperiodcentered\, caderno amarelo \,\textperiodcentered\, Q101}}
Mesmo para peixes de aqu\'ario, como o peixe arco-\'iris, a temperatura da
\'agua fora da faixa ideal (26 $^\circ$C a 28 $^\circ$C), bem como sua
varia\c{c}\~ao brusca, pode afetar a sa\'ude do animal. Para manter a
temperatura da \'agua dentro do aqu\'ario na m\'edia desejada, utilizam-se
dispositivos de aquecimento com termostato. Por exemplo, para um aqu\'ario
de 50 L, pode-se utilizar um sistema de aquecimento de 50 W otimizado
para suprir sua taxa de resfriamento. Essa taxa pode ser considerada
praticamente constante, j\'a que a temperatura externa ao aqu\'ario \'e
mantida pelas estufas. Utilize para a \'agua o calor espec\'ifico
4,0 kJ$\cdot$kg$^{-1}\cdot$K$^{-1}$ e a densidade 1 kg$\cdot$L$^{-1}$.

Se o sistema de aquecimento for desligado por 1 h, qual o valor mais
pr\'oximo para a redu\c{c}\~ao da temperatura da \'agua do aqu\'ario?
\begin{enumerate}
  \item 4,0 $^\circ$C
  \item 3,6 $^\circ$C
  \item 0,9 $^\circ$C
  \item 0,6 $^\circ$C
  \item 0,3 $^\circ$C
\end{enumerate}
\end{questao}


% ===================================================================
\parte{Trabalho e Energia}
% ===================================================================

\begin{questao}{6}{\fonte{ENEM 2007 \,\textperiodcentered\, prova amarela \,\textperiodcentered\, Q57}}
\figuraoriginal{0.8\linewidth}{q2007_057_mochila}

Com o projeto de mochila ilustrado acima, pretende-se aproveitar, na
gera\c{c}\~ao de energia el\'etrica para acionar dispositivos eletr\^onicos
port\'ateis, parte da energia desperdi\c{c}ada no ato de caminhar. As
transforma\c{c}\~oes de energia envolvidas na produ\c{c}\~ao de eletricidade
enquanto uma pessoa caminha com essa mochila podem ser assim
esquematizadas:

\figuraoriginal{0.5\linewidth}{q2007_057_esquema}

As energias I e II, representadas no esquema acima, podem ser
identificadas, respectivamente, como
\begin{enumerate}
  \item cin\'etica e el\'etrica.
  \item t\'ermica e cin\'etica.
  \item t\'ermica e el\'etrica.
  \item sonora e t\'ermica.
  \item radiante e el\'etrica.
\end{enumerate}
\end{questao}


% ===================================================================
\parte{Leis de Newton}
% ===================================================================

\begin{questao}{7}{\fonte{ENEM 2023 \,\textperiodcentered\, caderno amarelo \,\textperiodcentered\, Q130}}
Uma equipe de seguran\c{c}a do transporte de uma empresa avalia o
comportamento das tens\~oes que aparecem em duas cordas, 1 e 2, usadas para
prender uma carga de massa $M = 200$ kg na carroceria, conforme a
ilustra\c{c}\~ao. Quando o caminh\~ao parte do repouso, sua acelera\c{c}\~ao \'e
constante e igual a 3 m/s$^2$ e, quando ele \'e freado bruscamente, sua
frenagem \'e constante e igual a 5 m/s$^2$. Em ambas as situa\c{c}\~oes, a
carga encontra-se na imin\^encia de movimento, e o sentido do movimento do
caminh\~ao est\'a indicado na figura. O coeficiente de atrito est\'atico entre
a caixa e o assoalho da carroceria \'e igual a 0,2. Considere a acelera\c{c}\~ao
da gravidade igual a 10 m/s$^2$, as tens\~oes iniciais nas cordas iguais a
zero e as duas cordas ideais.

\figuraoriginal{0.62\linewidth}{q2023_130_caminhao}

Nas situa\c{c}\~oes de acelera\c{c}\~ao e frenagem do caminh\~ao, as tens\~oes nas
cordas 1 e 2, em newton, ser\~ao
\begin{enumerate}
  \item acelera\c{c}\~ao: $T_1 = 0$ e $T_2 = 200$; frenagem: $T_1 = 600$ e
        $T_2 = 0$.
  \item acelera\c{c}\~ao: $T_1 = 0$ e $T_2 = 200$; frenagem: $T_1 = 1\,400$ e
        $T_2 = 0$.
  \item acelera\c{c}\~ao: $T_1 = 0$ e $T_2 = 600$; frenagem: $T_1 = 600$ e
        $T_2 = 0$.
  \item acelera\c{c}\~ao: $T_1 = 560$ e $T_2 = 0$; frenagem: $T_1 = 0$ e
        $T_2 = 960$.
  \item acelera\c{c}\~ao: $T_1 = 640$ e $T_2 = 0$; frenagem: $T_1 = 0$ e
        $T_2 = 1\,040$.
\end{enumerate}
\end{questao}

% ===================================================================
\newpage
\parte{Gabarito}
% ===================================================================

\begin{center}
{\renewcommand{\arraystretch}{1.25}
\begin{tabular}{ccccccc}
  \toprule
  1 & 2 & 3 & 4 & 5 & 6 & 7 \\
  \midrule
  E & C & C & D & C & A & A \\
  \bottomrule
\end{tabular}}
\end{center}
\end{document}
